% =============================================================================
% Analisador espectral do Morphe: como funciona, do PC a FPGA (14/10/2026)
%
% Documento independente do manual, mas com o mesmo estilo (cores por
% camada: cliente azul, HPS verde, FPGA laranja). Figuras em TikZ, no
% proprio arquivo. Compila com pdfLaTeX (Overleaf: arquivo principal este).
%
% Fonte: branch estagio/v1.8-codec. Medidas: docs/DIARIO.md e
% docs/PRECISAO-NUMERICA.md.
% =============================================================================
\documentclass[11pt,a4paper,oneside]{article}

\usepackage[utf8]{inputenc}
\usepackage[T1]{fontenc}
\usepackage{lmodern}
\usepackage[brazilian]{babel}
\usepackage[a4paper,margin=2.5cm]{geometry}
\usepackage{microtype}
\usepackage{amsmath}
\setlength{\emergencystretch}{3em}
\usepackage{xcolor}
\usepackage{booktabs}
\usepackage{tabularx}
\usepackage{array}
\usepackage{listings}
\usepackage{caption}
\usepackage{float}
\usepackage{tikz}
\usetikzlibrary{positioning,arrows.meta,fit,backgrounds,calc,
                decorations.pathreplacing,shapes.geometric}
\usepackage[hidelinks]{hyperref}

% ---- estilo das figuras (o mesmo de docs/manual/figuras/estilo.tex) --------
\definecolor{cCliente}{HTML}{2F6DB5}
\definecolor{cHps}{HTML}{2E8B57}
\definecolor{cFpga}{HTML}{D2691E}
\definecolor{cRede}{HTML}{6B6B6B}
\definecolor{opFft}{HTML}{7D3C98}
\definecolor{opAdc}{HTML}{16877F}

\tikzset{
  >={Stealth[length=2.2mm]},
  every picture/.append style={font=\small},
  caixa/.style={draw, rounded corners=2pt, align=center, inner sep=4pt,
                minimum height=8mm},
  cliente/.style={caixa, draw=cCliente, fill=cCliente!10},
  hps/.style={caixa, draw=cHps, fill=cHps!10},
  fpga/.style={caixa, draw=cFpga, fill=cFpga!12},
  rede/.style={caixa, draw=cRede, fill=cRede!8},
  grupo/.style={draw, rounded corners=4pt, inner sep=6pt},
  seta/.style={->, thick},
  rotulo/.style={font=\scriptsize, align=center},
}

% ---- codigo ----------------------------------------------------------------
\lstset{
  basicstyle=\ttfamily\footnotesize,
  columns=fullflexible,
  keepspaces=true,
  frame=single,
  rulecolor=\color{black!25},
  backgroundcolor=\color{black!3},
  xleftmargin=4pt, xrightmargin=4pt,
  aboveskip=8pt, belowskip=8pt,
  showstringspaces=false,
  commentstyle=\color{black!55}\itshape,
  keywordstyle=\color{cCliente!80!black}\bfseries,
}
\lstdefinestyle{py}{language=Python}
\lstdefinestyle{c}{language=C, keywordstyle=\color{cHps!80!black}\bfseries}
\lstdefinestyle{txt}{language={}}

% Nome de arquivo ou comando, sempre em pe.
\newcommand{\arq}[1]{{\upshape\ttfamily\detokenize{#1}}}

\title{\textbf{Analisador espectral do Morphe}\\[4pt]
  \large Como funciona, do programa em Python ao hardware na FPGA}
\author{Antonio Nicassio Santos Lima\\
  \small Estágio no LabPS/UEFS --- supervisão: Prof.\ Armando S.\ Sanca}
\date{14 de outubro de 2026}

\begin{document}
\maketitle

\begin{abstract}
O analisador espectral do Morphe mostra, quadro a quadro, o espectro de um
sinal real ligado ao conversor analógico-digital da placa DE1-SoC. Cada quadro
atravessa as três camadas do sistema: o programa em Python, no computador,
pede a captura e prepara o sinal; o servidor em C, no processador ARM da placa,
comanda o hardware; e a FPGA amostra o sinal em instantes fixados por um
contador e calcula a FFT. Este documento segue um quadro do começo ao fim,
camada por camada, e explica cada conta e cada escala no caminho. Cada
afirmação aponta para o arquivo e a linha de onde saiu (branch
\arq{estagio/v1.8-codec}); as medidas vêm do diário do estágio.
\end{abstract}

\tableofcontents

% =============================================================================
\section{Visão geral}
\label{sec:visao}

A figura~\ref{fig:visao} mostra as peças envolvidas. O usuário abre a janela
``Analisador espectral (ADC $\to$ FFT na FPGA)'' na tela inicial do aplicativo
(\arq{Python/morphe_app.py:110}), escolhe canal, taxa de amostragem, número de
pontos e janela, e clica em \emph{Iniciar}. A partir daí, até clicar em
\emph{Parar}, o programa repete o ciclo da figura~\ref{fig:sequencia}.

\begin{figure}[htbp]
\centering
\begin{tikzpicture}[node distance=4mm, every node/.append style={font=\scriptsize}]
  % PC
  \node[cliente, minimum width=34mm] (jan) {\texttt{analisador\_window.py}\\janela, laço de quadros,\\contas do espectro};
  \node[cliente, minimum width=34mm, below=of jan] (aq) {\texttt{aquisicao.py}\\fs, palavra do ADC, volts};
  \node[cliente, minimum width=34mm, below=of aq] (bl) {\texttt{blocos.py}\\FFT de 1024 e FFT longa};
  \node[cliente, minimum width=34mm, below=of bl] (pr) {\texttt{morphe\_protocol.py}\\Q15.8, TCP};
  \node[grupo, draw=cCliente, fit=(jan)(aq)(bl)(pr), label={[cCliente]above:PC (Python)}] (gpc) {};

  % HPS
  \node[hps, minimum width=28mm, right=17mm of aq] (srv) {\texttt{morphe\_server.c}\\\texttt{handle\_adc}};
  \node[hps, minimum width=28mm, below=of srv] (srv2) {\texttt{morphe\_server.c}\\\texttt{handle\_fft}};
  \node[grupo, draw=cHps, fit=(srv)(srv2), label={[cHps]above:HPS (ARM, Linux)}] (ghps) {};

  % FPGA
  \node[fpga, minimum width=26mm, right=12mm of srv] (adc) {\texttt{adc\_captura.v}};
  \node[fpga, minimum width=26mm, right=12mm of srv2] (fft) {\texttt{fft\_wrapper.v}\\+ IP FFT};
  \node[fpga, minimum width=26mm, above=6mm of adc, fill=white] (ram1) {RAM \texttt{adc\_buf}};
  \node[fpga, minimum width=26mm, below=6mm of fft, fill=white] (ram2) {RAMs \texttt{fft\_xn}/\texttt{fft\_yn}};
  \node[grupo, draw=cFpga, fit=(adc)(fft)(ram1)(ram2), label={[cFpga]above:FPGA (Cyclone V)}] (gf) {};

  \node[caixa, draw=opAdc, fill=opAdc!10, right=5mm of gf.east |- adc] (ltc) {LTC2308\\(ADC)};
  \draw[<->, thick, opAdc] (adc) -- (ltc);
  \node[rotulo, above=1mm of ltc] {J15};

  \draw[seta, cRede] (gpc.east |- srv) -- node[above, rotulo]{OP\_ADC} (ghps.west |- srv);
  \draw[seta, cRede] (gpc.east |- srv2) -- node[above, rotulo]{OP\_FFT} (ghps.west |- srv2);
  \draw[<->, thick] (srv) -- node[above, rotulo]{PIOs} (adc);
  \draw[<->, thick] (srv2) -- node[above, rotulo]{PIOs} (fft);
  \draw[<->, thick] (adc) -- (ram1);
  \draw[<->, thick] (fft) -- (ram2);
\end{tikzpicture}
\caption{As peças do analisador, nas três camadas. As setas cinzas são
pedidos TCP; ``PIOs'' são os registradores de controle que o ARM escreve e lê
pela ponte leve.}
\label{fig:visao}
\end{figure}

Três ideias valem para tudo o que vem depois:

\begin{itemize}
\item \textbf{Captura e FFT são pedidos separados}, cada um numa conexão TCP
  própria. O servidor atende uma conexão por vez; entre um quadro e o próximo,
  outra pessoa pode usar a mesma placa.
\item \textbf{Quem calcula a FFT é a FPGA}, com o IP FFT da Intel (1024 pontos,
  ponto flutuante em bloco). O PC só prepara o sinal e interpreta o resultado.
  O NumPy refaz a mesma conta apenas para mostrar quanto as duas concordam.
\item \textbf{Acima de 1024 pontos} (4096 e 16\,384), o PC divide a FFT em
  várias FFTs de 1024 na placa e as junta pelo algoritmo de quatro passos. O
  resultado é exato, sem aproximação.
\end{itemize}

\begin{figure}[htbp]
\centering
\begin{tikzpicture}[
    ator/.style={caixa, minimum width=27mm, font=\scriptsize},
    msg/.style={->, thick},
    nota/.style={font=\scriptsize, align=left, fill=white, inner sep=1pt},
  ]
  \def\xa{0}\def\xb{34mm}\def\xc{72mm}\def\xd{112mm}
  \def\fim{-156mm}

  \node[ator, cliente] (a) at (\xa,0) {Janela\\(\texttt{analisador\_window})};
  \node[ator, cliente] (b) at (\xb,0) {\texttt{aquisicao}, \texttt{blocos},\\\texttt{morphe\_protocol}};
  \node[ator, hps]     (c) at (\xc,0) {\texttt{morphe\_server}};
  \node[ator, fpga]    (d) at (\xd,0) {\texttt{adc\_captura},\\\texttt{fft\_wrapper}};

  \foreach \n in {a,b,c,d}
    \draw[densely dashed, black!40] (\n.south) -- (\n.south |- 0,\fim);

  \fill[cHps!25] ($(\xc,-20mm)+(-1mm,0)$) rectangle ($(\xc,-56mm)+(1mm,0)$);
  \fill[cHps!25] ($(\xc,-96mm)+(-1mm,0)$) rectangle ($(\xc,-134mm)+(1mm,0)$);

  % captura
  \draw[msg] (\xa,-12mm) -- node[above, rotulo]{\texttt{capturar}(n, fs, canal)} (\xb,-12mm);
  \draw[msg, cRede] (\xb,-20mm) -- node[above, rotulo]{OP\_ADC: n, divisor, config} (\xc,-20mm);
  \draw[msg] (\xc,-28mm) -- node[above, rotulo]{divisor, n, config; \emph{start} = 1} (\xd,-28mm);
  \node[nota, right] at ($(\xd,-36mm)+(1.5mm,0)$) {$n+1$ conversões,\\uma a cada\\divisor ciclos};
  \draw[msg, dashed] (\xd,-44mm) -- node[above, rotulo]{\emph{done} = 1} (\xc,-44mm);
  \draw[msg, dashed] (\xd,-50mm) -- node[above, rotulo]{lê a RAM \texttt{adc\_buf}} (\xc,-50mm);
  \draw[msg, cRede] (\xc,-56mm) -- node[above, rotulo]{$n$ códigos \texttt{int32}} (\xb,-56mm);
  \draw[msg] (\xb,-62mm) -- node[above, rotulo]{volts (1 código = 1\,mV)} (\xa,-62mm);

  % preparo
  \draw[msg] (\xa,-68mm) arc[start angle=90, end angle=-90, radius=2.5mm]
             node[right, nota, pos=0.5, xshift=1mm] {tira o DC, aplica a janela};
  \draw[msg] (\xa,-80mm) -- node[above, rotulo]{FFT de $x_w[n]$} (\xb,-80mm);
  \draw[msg] (\xb,-84mm) arc[start angle=90, end angle=-90, radius=2.5mm]
             node[right, nota, pos=0.5, xshift=1mm] {escala para Q15.8};

  % FFT
  \draw[msg, cRede] (\xb,-96mm) -- node[above, rotulo]{OP\_FFT: 1024 $\times$ Q15.8} (\xc,-96mm);
  \draw[msg] (\xc,-104mm) -- node[above, rotulo]{escreve \texttt{fft\_xn}; \emph{start} = 1} (\xd,-104mm);
  \node[nota, right] at ($(\xd,-110mm)+(1.5mm,0)$) {IP FFT\\(1024, BFP)};
  \draw[msg, dashed] (\xd,-116mm) -- node[above, rotulo]{\emph{done} = 1} (\xc,-116mm);
  \draw[msg, dashed] (\xd,-122mm) -- node[above, rotulo]{lê \texttt{fft\_yn} e o expoente} (\xc,-122mm);
  \draw[msg] (\xc,-125mm) arc[start angle=90, end angle=-90, radius=2mm]
             node[right, nota, pos=0.5, xshift=1mm] {desfaz BFP e Q15.8};
  \draw[msg, cRede] (\xc,-134mm) -- node[above, rotulo]{1024 complexos \texttt{float32}} (\xb,-134mm);
  \draw[msg] (\xb,-140mm) -- node[above, rotulo]{$X[k]$ (escala desfeita)} (\xa,-140mm);
  \draw[msg] (\xa,-146mm) arc[start angle=90, end angle=-90, radius=2.5mm]
             node[right, nota, pos=0.5, xshift=1mm] {volts de pico, pico, compara com NumPy; desenha};
\end{tikzpicture}
\caption{Um quadro do analisador com 1024 pontos. Com 4096 ou 16\,384, o
trecho da FFT se repete 4 ou 16 vezes (seção~\ref{sec:longa}).}
\label{fig:sequencia}
\end{figure}

% =============================================================================
\section{O programa em Python}
\label{sec:python}

\subsection{Onde está cada coisa}

\begin{table}[H]
\centering
\small
\begin{tabularx}{\linewidth}{@{}l X@{}}
\toprule
arquivo & papel no analisador \\
\midrule
\arq{analisador_window.py} & a janela: opções, laço de quadros, contas do espectro, gráficos \\
\arq{aquisicao.py} & captura pelo ADC: palavra de configuração, fs real, volts, saturação \\
\arq{blocos.py} & \texttt{FftPlaca} (uma FFT na placa) e \texttt{fft\_longa} (quatro passos) \\
\arq{morphe_protocol.py} & monta e decodifica as mensagens binárias; \texttt{TcpClient} \\
\arq{dsp_core.py} & codificação Q15.8 (\texttt{fft\_q1508\_encode}), \texttt{.mrph} \\
\arq{morphe_config.py} & os mesmos limites de \arq{C/morphe_config.h} (1024, 32\,768, 50\,MHz) \\
\bottomrule
\end{tabularx}
\caption{Os módulos Python do analisador (pasta \arq{Python/}).}
\end{table}

\subsection{As opções da janela}

Montadas em \texttt{AnalisadorWindow.\_build\_sidebar}
(\arq{analisador_window.py:225}):

\begin{table}[H]
\centering
\small
\begin{tabularx}{\linewidth}{@{}l X@{}}
\toprule
opção & o que faz \\
\midrule
Modo / Canal & simples (CH0 a CH7 contra o terra, 0 a 4,095\,V) ou diferencial
  (CH0$-$CH1 \dots{} CH6$-$CH7, $-2{,}048$ a $+2{,}047$\,V) \\
fs (Hz) & taxa de amostragem pedida, de 1\,kHz a 200\,kHz \\
Pontos & 1024 (1 FFT na FPGA), 4096 (4 FFTs) ou 16\,384 (16 FFTs) \\
Janela & Hann, Blackman-Harris, flat-top, retangular \\
Média & nenhuma, média de potência de 4 ou 16 quadros, retenção de pico \\
Tirar o nível DC & subtrai a média do quadro antes da FFT \\
Mostrar a fase & terceiro gráfico, $\angle X(f)$ \\
Stem & desenha as amostras como hastes, sem ligar os pontos \\
\bottomrule
\end{tabularx}
\caption{Opções da janela do analisador.}
\end{table}

O texto de informação (\texttt{\_atualiza\_info}, linha 342) mostra a
\textbf{fs real}, a resolução em frequência $f_s/N$ e a duração do quadro
$N/f_s$, porque a fs pedida raramente é a que o hardware usa
(seção~\ref{sec:captura-py}).

\subsection{O laço de quadros}

\texttt{\_on\_iniciar} (\arq{analisador_window.py:368}) cria o cliente TCP a
partir do painel de conexão e dispara uma \textbf{thread de trabalho}
(\texttt{worker}, linha 395), que repete até o evento \texttt{parar} ser
acionado:

\begin{lstlisting}[style=py]
cap = aq.capturar(client, n, fs, modo, canal, guardar=False)   # pedido 1: ADC
q = processar_quadro(cap.volts, cap.fs, janela, tirar_dc, fft)  # pedido(s) 2: FFT
with self._trava_q:
    self._novo = q                                              # entrega a tela
\end{lstlisting}

A thread não desenha: o Tk só pode ser usado pela thread principal. Por isso
há dois ritmos independentes:

\begin{itemize}
\item a thread produz quadros tão rápido quanto a placa e a rede permitem, e
  deixa sempre só o \textbf{último} em \texttt{self.\_novo}, protegido por uma
  trava;
\item a thread do Tk confere a cada \texttt{TELA\_MS} = 50\,ms
  (\texttt{\_tela}, linha 504) se há quadro novo e, se houver, acumula a média
  e redesenha.
\end{itemize}

Se a tela for mais lenta que a placa, quadros intermediários deixam de ser
desenhados, mas a análise não atrasa. As opções que a thread lê (janela, tirar
DC) ficam no dicionário \texttt{self.\_cfg}, copiado pela thread do Tk em
\texttt{\_zera\_media} (linha 359): variáveis do Tk não podem ser lidas de
outra thread.

\subsection{A captura}
\label{sec:captura-py}

\texttt{aquisicao.capturar} (\arq{aquisicao.py:139}) faz três contas antes de
qualquer byte sair do PC.

\paragraph{A fs que o hardware realmente usa.} A FPGA conta ciclos de um
relógio de 50\,MHz, e o período de amostragem tem de ser um número inteiro de
ciclos, o \emph{divisor}:
\[
  \text{divisor} = \operatorname{round}\!\left(\frac{50\,\text{MHz}}{f_s}\right)
  \in [250,\ 50\,000],
  \qquad
  f_{s,\text{real}} = \frac{50\,\text{MHz}}{\text{divisor}} .
\]
Pedir 44\,100\,Hz dá divisor 1134 e $f_{s,\text{real}} = 44\,091{,}7$\,Hz
(\texttt{divisor\_para} e \texttt{fs\_de}, linhas 86 e 93). \textbf{Todas as
contas seguintes usam a fs real}; é ela que aparece no eixo de frequência.

\paragraph{A palavra de configuração do LTC2308.} \texttt{palavra\_config}
(linha 63) monta os 6 bits \texttt{S/D O/S S1 S0 UNI SLP}, que escolhem o canal,
o modo simples ou diferencial e a saída unipolar ou bipolar. No modo simples, o
canal é codificado como $\{S1, S0, O/S\}$ (tabela~1 da folha de dados).
\texttt{SLP} fica sempre em 0: dormir desliga a referência interna, que leva
200\,ms para voltar.

\paragraph{O tempo de espera.} A captura leva $(n+1)\cdot
\text{divisor}/50\,\text{MHz}$ segundos (o primeiro quadro é descartado,
seção~\ref{sec:adc-hw}). O cliente usa um timeout de pelo menos essa duração
mais 10\,s: 32\,768 amostras a 1\,kHz levam 33\,s, e o padrão do painel TCP é
10\,s.

A resposta são \textbf{códigos inteiros} do conversor: 0 a 4095 no modo simples
e $-2048$ a 2047 no diferencial (o servidor já estendeu o sinal). Com a
referência interna de 4,096\,V, \textbf{1 código = 1\,mV}, e a conversão para
volts é \texttt{codigos * 0.001}. O cliente confere que voltaram exatamente $n$
amostras e que o divisor usado (campo \texttt{extra} da resposta) é o pedido.
\texttt{saturou} (linha 357) avisa se algum código bateu no limite da faixa.

\subsection{O processamento de um quadro}
\label{sec:quadro}

\texttt{processar\_quadro} (\arq{analisador_window.py:120}) é uma função pura,
sem Tk: recebe os volts e uma função \texttt{fft}, e devolve um dicionário com
tudo o que a tela mostra. Os testes de bancada a chamam com
\texttt{np.fft.fft} no lugar da placa.

\paragraph{Passo 1: nível DC e janela.}
\[
  \bar{x} = \frac{1}{N}\sum_{n} x[n], \qquad
  x_w[n] = \bigl(x[n] - \bar{x}\bigr)\, w[n] .
\]
No modo simples o sinal nunca é negativo: um gerador com offset de 2\,V produz
uma raia enorme em $f=0$ que, vazando pela janela, esconderia os tons baixos.
Subtrair a média remove essa raia.

As janelas (\texttt{janela}, linha 109) são \textbf{periódicas}, de soma
conhecida: Hann, Blackman-Harris de 4 termos, flat-top de 5 termos e
retangular. A de Hann, por exemplo, é \texttt{np.hanning(n+1)} sem a última
amostra. A escolha é um compromisso
(tabela~\ref{tab:janelas}).

\begin{table}[H]
\centering
\small
\begin{tabularx}{\linewidth}{@{}l l X@{}}
\toprule
janela & erro máximo de amplitude & uso \\
\midrule
Retangular & até 3,9\,dB a menos & só com número inteiro de períodos no quadro \\
Hann & até 1,4\,dB & uso geral \\
Blackman-Harris & até 0,8\,dB; lóbulos abaixo de $-92$\,dB & separar tons fracos de fortes \\
Flat-top & menos de 0,02\,dB & \textbf{medir amplitude}; pico mais largo \\
\bottomrule
\end{tabularx}
\caption{Erro de amplitude de um tom que cai entre dois bins, por janela
(comentário de \arq{analisador_window.py:48}, conferido em bancada em 30/09).}
\label{tab:janelas}
\end{table}

\paragraph{Passo 2: a FFT, na FPGA.} $X = \texttt{fft}(x_w)$, com o tempo
medido. A função \texttt{fft} é montada em \texttt{\_on\_iniciar} (linha 389):

\begin{lstlisting}[style=py]
fft_placa = blocos.FftPlaca(client, normalizar=True)
def fft(x):
    return blocos.fft_longa(x, fft_placa) if x.size > 1024 else fft_placa(x)
\end{lstlisting}

Os detalhes estão nas seções \ref{sec:fftplaca} e \ref{sec:longa}.

\paragraph{Passo 3: de $|X[k]|$ para volts de pico.} Uma senoide de amplitude
$A$, janelada por $w$, produz no seu bin $|X[k]| = A \cdot \tfrac{1}{2}\sum_n
w[n]$. Então
\[
  \hat{A}[k] = \frac{|X[k]|}{\tfrac{1}{2}\sum_n w[n]}, \qquad
  \hat{A}[0] = \frac{|X[0]|}{\sum_n w[n]} ,
\]
porque no DC o ganho é a soma inteira, e não a metade. O espectro mostrado é
\textbf{completo, de 0 a $f_s$} (os $N$ bins), com uma linha pontilhada em
$f_s/2$: como o sinal é real, $X[N-k] = X^*[k]$, e um tom em $f$ aparece também
em $f_s - f$, com a mesma amplitude. O eixo vertical é \textbf{linear, em volts
de pico}, conforme decidido em 01/10 (\arq{docs/DECISOES.md}).

\paragraph{Passo 4: o pico.} É procurado só na metade de baixo, ignorando os
três primeiros bins (\texttt{BINS\_DC}, o que sobra do DC depois da janela). A
frequência é refinada por uma \textbf{parábola} nos três bins em volta do
máximo, em dB (linha 156). Com $a, b, c$ os valores em dB dos bins $k-1$, $k$ e
$k+1$:
\[
  \delta = \frac{1}{2}\,\frac{a-c}{a-2b+c}, \qquad
  f_{\text{pico}} = (k + \delta)\,\frac{f_s}{N}, \qquad |\delta| \le 0{,}5 .
\]
O lóbulo de uma janela é quase parabólico em escala logarítmica, e isso
corrige o erro de até meio bin do índice inteiro. O limite em $\pm 0{,}5$ bin
existe porque, perto do DC ou só com ruído, a parábola pode apontar longe.

\paragraph{Passo 5: a prova.} O NumPy calcula \texttt{np.fft.fft} do
\textbf{mesmo} $x_w$ enviado à placa, e a tela mostra
\[
  \text{SNR} = 20\log_{10}
  \frac{\lVert X_{\text{NumPy}} \rVert}{\lVert X_{\text{FPGA}} - X_{\text{NumPy}} \rVert}
  \quad\text{(dB)}
\]
como ``FPGA $\times$ NumPy''. Medido na placa~1 em 30/09/2026: \textbf{97,6\,dB}
em todos os quadros com 1024 pontos, e 97,4 a 97,8\,dB com 4096. É a
confirmação, a cada quadro, de que o espectro saiu mesmo da FPGA e com que
precisão.

\subsection{Uma FFT de 1024 pontos na placa}
\label{sec:fftplaca}

\texttt{FftPlaca} (\arq{blocos.py:106}):

\begin{lstlisting}[style=py]
x = dsp.pad_zeros_to(x, 1024)
req, escala = build_fft_request(x, normalizar=True)
return decode_fft_response(self.client.request(req), escala)
\end{lstlisting}

O hardware recebe números em \textbf{Q15.8}: inteiros de 24 bits com sinal,
dos quais 8 são fracionários. A faixa é $\pm 32\,768$ e o passo, $1/256$. Um
sinal do ADC tem no máximo $\pm 2$\,V; enviado como está, usaria uns 10 dos 24
bits, e o erro de arredondamento seria grande em relação ao sinal.

Por isso a \textbf{normalização} (\texttt{escala\_para\_q1508},
\arq{morphe_protocol.py:198}):
\[
  s = \frac{0{,}98 \cdot 32\,767{,}996}{\max_n |x_w[n]|}, \qquad
  q[n] = \operatorname{round}\bigl(256 \cdot s\, x_w[n]\bigr) ,
\]
com $q$ saturado em 24 bits pela função \texttt{fft\_q1508\_encode}. Na volta,
\texttt{decode\_fft\_response} (linha 450) divide: $X = X_{\text{placa}}/s$. A conta é
exata, porque a DFT é linear: $\operatorname{DFT}(s\,x)/s =
\operatorname{DFT}(x)$. Muda só o erro de quantização, que cai 6,02\,dB por bit
de faixa recuperado, até o teto de cerca de 92\,dB do IP. Cada bloco é
normalizado por conta própria. As medidas estão na seção~3 de
\arq{docs/PRECISAO-NUMERICA.md}.

Cada amostra Q15.8 viaja num \texttt{int32} big-endian, com os 8 bits de cima
repetindo o bit de sinal; o hardware só usa os 24 de baixo.

\subsection{FFT de 4096 e 16\,384 pontos}
\label{sec:longa}

\texttt{fft\_longa} (\arq{blocos.py:241}) calcula uma DFT de $N = 1024\,M$
pontos pelo algoritmo de \textbf{quatro passos} (Cooley-Tukey de raiz $M$).
Escrevendo $n = n_1 + M n_2$ e $k = k_2 + 1024\,k_1$, com $0 \le n_1, k_1 < M$
e $0 \le n_2, k_2 < 1024$:
\[
  X[k_2 + 1024\,k_1] =
  \sum_{n_1=0}^{M-1}
  \underbrace{e^{-j2\pi n_1 k_1/M}}_{\text{passo 4: DFT de } M}
  \;
  \underbrace{e^{-j2\pi n_1 k_2/N}}_{\text{passo 3: giro}}
  \;
  \underbrace{\sum_{n_2=0}^{1023} x[n_1 + M n_2]\, e^{-j2\pi n_2 k_2/1024}}_{\text{passos 1 e 2: } A[n_1][k_2]\text{, na placa}} .
\]

\begin{enumerate}
\item separa $x$ em $M$ subsequências decimadas, $x_{n_1}[n_2] = x[n_1 + M n_2]$;
\item calcula a FFT de 1024 de cada uma \textbf{na placa}: $M$ pedidos OP\_FFT;
\item multiplica pelos fatores de giro $e^{-j2\pi n_1 k_2/N}$;
\item calcula, no PC, $1024$ DFTs de $M$ pontos ao longo de $n_1$.
\end{enumerate}

No código, os passos 3 e 4 são duas linhas:

\begin{lstlisting}[style=py]
giro = np.exp(-2j*np.pi * np.outer(np.arange(M), np.arange(1024)) / N)
return np.fft.fft(A * giro, axis=0).reshape(-1)
\end{lstlisting}

O trabalho pesado ($M$ FFTs de 1024) fica com a FPGA; o PC faz só $1024\,M$
multiplicações complexas e as DFTs minúsculas de $M$ pontos. Com 16\,384 pontos
são 16 pedidos à placa por quadro, o que reduz a taxa de quadros.

\subsection{A tela: média, fase, salvar}

\begin{itemize}
\item \textbf{Média} (\texttt{\_acumula}, linha 488): sempre sobre a
  \textbf{potência} $\hat{A}^2$, nunca sobre a amplitude ou o complexo. Média de
  $K$ quadros: $\sqrt{\operatorname{média}(\hat{A}^2)}$ dos últimos $K$;
  retenção de pico: $\sqrt{\max \hat{A}^2}$ desde o começo. Mudar a janela, o
  DC ou o tipo de média zera o acumulador.
\item \textbf{Fase} (\texttt{fase\_visivel}, linha 99): $\angle X[k]$ só nos
  bins com amplitude de pelo menos 2\,\% da maior. Nos outros, $X[k]$ é ruído e
  a fase é um número sorteado entre $-\pi$ e $\pi$; por isso aparecem como
  lacuna. A fase é sempre a do quadro atual: cada captura começa num ponto
  diferente da onda, e uma média de fase não teria sentido.
\item \textbf{Medidas do quadro}: pico (Hz e mV), nível DC, tempo de captura,
  tempo da FFT, quadros por segundo, FPGA $\times$ NumPy e aviso de saturação.
\item \textbf{Salvar quadro} (\texttt{\_on\_salvar}, linha 561): grava em
  \texttt{.mrph} o $x_w$ (exatamente o que foi para a FPGA) e o $X$ da placa,
  no formato que o Comparador abre.
\end{itemize}

% =============================================================================
\section{O protocolo entre o PC e a placa}
\label{sec:protocolo}

Definido em dois arquivos que espelham um ao outro, \arq{Python/morphe_protocol.py}
e \arq{C/morphe_protocol.h}. TCP na porta 5000, \textbf{uma conexão por
pedido}, todos os inteiros em big-endian.

\begin{table}[H]
\centering
\small
\begin{tabularx}{\linewidth}{@{}l l X@{}}
\toprule
campo & tipo & valor \\
\midrule
magic & \texttt{u32} & \texttt{0x4D52504D} (``MRPM'') \\
version & \texttt{u16} & 1 \\
opcode & \texttt{u16} & 2 = FFT, 7 = ADC (há também CONV, FIR, IFFT, IIR, PING, ADC contínuo) \\
dtype & \texttt{u16} & 1 = \texttt{int32} \\
flags & \texttt{u16} & ADC: palavra de 6 bits do LTC2308; FFT: 0 \\
\texttt{n\_x} & \texttt{u32} & ADC: número de amostras; FFT: 1024 \\
\texttt{n\_h} & \texttt{u32} & ADC: divisor; FFT: 0 \\
\bottomrule
\end{tabularx}
\caption{Cabeçalho de pedido, 20 bytes (\texttt{struct.pack(">IHHHHII", ...)}).}
\end{table}

O cabeçalho da resposta também tem 20 bytes: magic \texttt{0x4D52504E}
(``MRPN''), version, opcode (espelhado), dtype, status (0 = OK),
\texttt{n\_out} e \texttt{extra}.

\begin{table}[H]
\centering
\small
\begin{tabularx}{\linewidth}{@{}l X X@{}}
\toprule
 & pedido & resposta \\
\midrule
OP\_ADC (7) & só o cabeçalho & $n$ códigos \texttt{int32}; \texttt{extra} = divisor usado \\
OP\_FFT (2) & cabeçalho + 1024 \texttt{int32} Q15.8 (4\,KiB) & 1024 complexos \texttt{float32}, re e im intercalados (8\,KiB) \\
\bottomrule
\end{tabularx}
\caption{Os dois pedidos do analisador.}
\end{table}

Em erro, o status é diferente de zero e a carga é a mensagem em texto
(\texttt{n\_out} = número de bytes). O Python a transforma em
\texttt{RuntimeError("erro do servidor: ...")}, que a janela mostra numa
caixa de diálogo. Por exemplo, com a FPGA ainda no bitstream de fábrica, a
resposta é o status 8, ``FFT: FPGA nao preparada -- rode ./morphe-up.sh nesta
placa''.

\texttt{TcpClient.request} (\arq{morphe_protocol.py:370}) abre a conexão,
desliga o algoritmo de Nagle (\texttt{TCP\_NODELAY}, que atrasaria mensagens
curtas), envia o pedido, lê exatamente os 20 bytes do cabeçalho, calcula o
tamanho da carga pelo opcode (FFT: $8\,n_\text{out}$; ADC: $4\,n_\text{out}$) e
lê exatamente esse tanto.

% =============================================================================
\section{O servidor em C}
\label{sec:servidor}

\arq{C/morphe_server.c} é um único processo, no Linux do ARM da placa,
executado como root por causa do \texttt{/dev/mem}.

\subsection{Como o ARM enxerga a FPGA}

O ARM e a FPGA estão no mesmo chip (Cyclone~V SoC), ligados por pontes AXI. Do
lado do ARM, cada memória e cada registrador da FPGA é um \textbf{endereço
físico}. \texttt{fpga\_init} (linha 413) abre \texttt{/dev/mem} e mapeia duas
janelas no espaço do processo com \texttt{mmap} (tabela~\ref{tab:mapa}).

\begin{table}[H]
\centering
\small
\begin{tabularx}{\linewidth}{@{}l l X@{}}
\toprule
janela & base física & o que há dentro (offsets de \arq{C/hps_0.h}) \\
\midrule
ponte HPS$\to$FPGA & \texttt{0xC0000000} &
  RAMs: \texttt{fft\_yn\_re} 0x18000, \texttt{fft\_yn\_imag} 0x19000,
  \texttt{fft\_xn\_re} 0x1A000, \texttt{fft\_xn\_imag} 0x1B000 (4\,KiB cada,
  1024 \texttt{int32}); \texttt{adc\_buf} 0x20000 (128\,KiB, 32\,768 \texttt{int32}) \\
ponte leve & \texttt{0xFF200000} &
  PIOs da FFT: \texttt{bfp\_exponent} 0x50, \texttt{done} 0x60,
  \texttt{inverse} 0x70, \texttt{start} 0x80. PIOs do ADC: \texttt{start}
  0xD0, \texttt{done} 0xE0, \texttt{config} 0xF0, \texttt{divisor} 0x100,
  \texttt{namostras} 0x110, \texttt{contador} 0x120 \\
\bottomrule
\end{tabularx}
\caption{As duas janelas de memória que o servidor mapeia.}
\label{tab:mapa}
\end{table}

Depois do mapeamento, escrever num registrador da FPGA é uma atribuição de
ponteiro: \texttt{*g\_pio\_adc\_divisor = divisor;}. Os ponteiros são
\texttt{volatile}, para o compilador não eliminar nem reordenar os acessos.
Os endereços de \arq{hps_0.h} saem do projeto do Platform Designer
(\arq{Quartus/gen_hps_header.py}): \textbf{se o bitstream muda de mapa, o
servidor tem de ser recompilado com o \arq{hps_0.h} novo.}

\subsection{Proteções antes de tocar na FPGA}

Acessar pela ponte um endereço que a FPGA não implementa \textbf{trava o
barramento do ARM}, sem timeout: a placa some da rede. Isso aconteceu em 21/09,
quando o servidor subia no boot com a FPGA ainda no bitstream de fábrica. Daí
duas proteções:

\begin{itemize}
\item \texttt{fpga\_preparada} (linha 495): só deixa uma operação seguir se
  existe o arquivo \texttt{/var/run/morphe-fpga-preparada}, que o
  \texttt{morphe-up.sh} cria depois de programar o \texttt{.sof} do Morphe. Na
  primeira vez, zera os PIOs de \emph{start}.
\item \texttt{adc\_recusa} (linha 1097), só para o ADC: compara o
  \emph{timestamp} do \texttt{sysid} gravado no bitstream com o
  \texttt{SYSID\_QSYS\_TIMESTAMP} com que o servidor foi compilado. Se forem
  diferentes, o bitstream é outro e os PIOs do ADC podem nem existir. Confere
  também o divisor (250 a 50\,000), o dtype e o bit SLP.
\end{itemize}

\subsection{O laço principal}

\texttt{main} (linha 1566) abre o socket na porta 5000 e repete
\texttt{accept} $\to$ \texttt{serve\_connection} $\to$ \texttt{close}.
\textbf{Uma conexão por vez}: enquanto uma captura ou FFT roda, a próxima
conexão espera na fila do \texttt{listen}. É isso que permite a várias pessoas
usarem a mesma placa: cada quadro do analisador são duas ou mais conexões
curtas, e as dos outros entram entre elas.

\texttt{serve\_connection} (linha 1524) lê o cabeçalho, confere magic e versão
e despacha pelo opcode. Em paralelo, uma thread responde na porta UDP 5001 quem
está usando a placa e qual operação está rodando (\texttt{thread\_estado}); é o
que a tela de descoberta mostra como ``ocupada''.

\subsection{OP\_ADC: \texttt{handle\_adc}}

Linha 1138, resumida (comentários sem acento, como no código):

\begin{lstlisting}[style=c]
*g_pio_adc_divisor   = divisor;      /* periodo em ciclos de 50 MHz */
*g_pio_adc_namostras = n;
*g_pio_adc_config    = cfg;          /* bit 6 = 0: captura unica */
*g_pio_adc_start = 0; usleep(1); *g_pio_adc_start = 1;  /* borda de subida */

wait_done(g_pio_adc_done, duracao_ms + 2000);           /* polling do done */
*g_pio_adc_start = 0;                                    /* libera o hardware */

for (i = 0; i < n; i++)                                  /* le a RAM pela ponte */
    tx[i] = adc_codigo(g_adc_buf[i], bipolar);           /* 12 bits -> int32 */
\end{lstlisting}

\begin{itemize}
\item O \textbf{aperto de mão} \emph{start}/\emph{done} é o mesmo em todos os
  blocos do Morphe: o ARM sobe \emph{start}; o hardware trabalha, sobe
  \emph{done} e o mantém em 1 até o ARM baixar \emph{start}. Assim nenhum dos
  dois lados perde o evento do outro.
\item \texttt{wait\_done} (linha 529) lê o PIO a cada 50\,$\mu$s, com timeout.
\item \texttt{adc\_codigo} (linha 1132) pega os 12 bits de baixo da palavra e,
  no modo bipolar (\texttt{UNI} = 0), estende o sinal do complemento de dois.
\item A resposta leva o \textbf{divisor no campo \texttt{extra}}, para o
  cliente confirmar a fs.
\end{itemize}

\subsection{OP\_FFT: \texttt{handle\_fft} e \texttt{fft\_run\_block}}
\label{sec:fft-srv}

Linhas 1375 e 1334, resumidas:

\begin{lstlisting}[style=c]
recv_exact(sock, rx_buf, 1024*4);              /* 1024 int32 Q15.8 */
decode_samples_to_q1508(rx_buf, 1024, x);      /* big-endian; satura em 24 bits */
for (i = 0; i < 1024; i++) {
    g_fft_xn_re[i] = x[i]; g_fft_xn_imag[i] = 0;   /* entrada real */
}
*g_pio_fft_inverse = 0;                        /* ANTES do start */
*g_pio_fft_start = 0; usleep(1); *g_pio_fft_start = 1;
wait_done(g_pio_fft_done, 5000);
*g_pio_fft_start = 0;

exp = *g_pio_fft_bfp_exponent & 0x3F;          /* 6 bits, com sinal */
for (k = 0; k < 1024; k++) {
    re[k] = g_fft_yn_re[k]; im[k] = g_fft_yn_imag[k];
}
escala = ldexpf(1.0f, -exp) / 256;             /* desfaz BFP e Q15.8 */
/* envia (float) re[k] * escala e (float) im[k] * escala, big-endian */
\end{lstlisting}

O ponto central é o \textbf{expoente de ponto flutuante em bloco} (BFP). Para
não estourar os 24 bits ao longo dos estágios, o IP divide os dados por 2
sempre que um estágio cresceria demais, e conta quantas vezes fez isso. Esse
número sai junto com o resultado: um único expoente para o bloco inteiro de
1024 bins. O servidor o lê do PIO \texttt{fft\_bfp\_exponent}, multiplica a
saída por $2^{-\text{exp}}$ para recuperar a escala verdadeira e divide por 256
para tirar o Q15.8 da entrada. O que chega ao Python é a DFT de $s\,x_w$ em
ponto flutuante; o Python divide por $s$ (seção~\ref{sec:fftplaca}).

O bit \texttt{inverse} é escrito \emph{antes} de subir o \emph{start} porque o
wrapper o trava na borda de subida (seção~\ref{sec:wrapper}). Cada FFT também
grava um arquivo \texttt{.mrph} de depuração no cartão da placa
(\texttt{save\_debug\_bundle\_fft}); os mais antigos são apagados
(\texttt{MORPHE\_DUMP\_MAX}, padrão 100).

% =============================================================================
\section{O hardware na FPGA}
\label{sec:fpga}

O projeto Quartus está em \arq{Quartus/}. O topo, \arq{ghrd_top.v}, instancia o
sistema do Platform Designer (\texttt{soc\_system}: HPS, pontes, RAMs
\emph{on-chip}, PIOs) e, ao lado dele, os blocos próprios do Morphe. Todos usam
o relógio \texttt{CLOCK\_50}, de 50\,MHz.

\subsection{As memórias de duas portas}

Cada RAM do caminho (\texttt{fft\_xn\_re}, \texttt{fft\_xn\_imag},
\texttt{fft\_yn\_re}, \texttt{fft\_yn\_imag}, \texttt{adc\_buf}) tem
\textbf{duas portas}:

\begin{itemize}
\item a \textbf{porta A}, no barramento, que o ARM lê e escreve pela ponte;
\item a \textbf{porta B}, exportada para o \arq{ghrd_top.v} e ligada
  diretamente ao bloco de hardware.
\end{itemize}

É por elas que os dados passam entre o software e o hardware, sem DMA: o ARM
escreve a entrada pela porta A, o hardware a lê pela porta B, escreve o
resultado pela porta B de outra RAM, e o ARM o lê pela porta A.

\subsection{O controlador do ADC: \texttt{adc\_captura.v}}
\label{sec:adc-hw}

Instanciado em \arq{ghrd_top.v:604}, com os pinos \texttt{ADC\_CONVST},
\texttt{ADC\_SCLK}, \texttt{ADC\_DIN} e \texttt{ADC\_DOUT} ligados ao LTC2308
e a porta B da RAM \texttt{adc\_buf} ($2^{15}$ palavras).

\paragraph{Por que um controlador próprio.} O IP do University Program converte
sem parar, no ritmo dele, e entrega só o último valor de cada canal. Ler esse
valor a cada $1/f_s$ deixa o instante da amostra incerto em alguns
microssegundos, e essa incerteza vira ruído que cresce com a frequência do
sinal (\arq{docs/ADC.md}). Aqui, \textbf{cada conversão começa exatamente no
ciclo em que o contador de período volta a zero}: o instante de amostragem é
fixado pelo hardware, com a incerteza de um ciclo de 50\,MHz (20\,ns). Para um
analisador espectral isso é essencial, porque a DFT supõe amostras igualmente
espaçadas.

\paragraph{A máquina de estados.}
\begin{description}
\item[\texttt{S\_IDLE}] na borda de subida de \emph{start}, trava
  \texttt{divisor} (no mínimo 250), \texttt{n\_amostras}, \texttt{config\_adc}
  e \texttt{continuo}, e zera o índice e o contador;
\item[\texttt{S\_PREP}] espera 16 ciclos, para garantir o tempo de aquisição
  do conversor;
\item[\texttt{S\_CAPT}] um \textbf{quadro por amostra}, contado pela variável
  \texttt{fase}, de 0 a $\text{divisor}-1$ (figura~\ref{fig:quadro-adc});
\item[\texttt{S\_DONE}] sobe \emph{done} e espera o ARM baixar \emph{start}.
\end{description}

\begin{figure}[htbp]
\centering
\begin{tikzpicture}[x=1mm, y=1mm, font=\scriptsize]
  % eixo de fase (fora de escala): 0, 96, 98, 146, divisor-1
  \def\yA{0}\def\yB{-14}\def\yC{-25}\def\yD{-35}
  \node[left] at (-2,\yA+2) {CONVST};
  \node[left] at (-2,\yB+2) {SCLK};
  \node[left] at (-2,\yC+2) {SDI / SDO};
  \node[left] at (-2,\yD+2) {grava na RAM};

  % CONVST
  \draw[thick, opAdc] (0,\yA) -- (0,\yA+4) -- (40,\yA+4) -- (40,\yA) -- (130,\yA) -- (130,\yA+4) -- (136,\yA+4);
  % SCLK: 12 pulsos entre 44 e 92
  \draw[thick, opAdc] (0,\yB) -- (44,\yB)
    \foreach \i in {0,...,11} { -- ++(2,0) -- ++(0,4) -- ++(2,0) -- ++(0,-4) }
    -- (136,\yB);
  % SDI/SDO: barramento de 12 bits
  \draw[thick, opAdc] (0,\yC+2) -- (44,\yC+2);
  \draw[thick, opAdc] (44,\yC+2) -- (45,\yC+4) -- (91,\yC+4) -- (92,\yC+2) -- (136,\yC+2);
  \draw[thick, opAdc] (44,\yC+2) -- (45,\yC) -- (91,\yC) -- (92,\yC+2);
  \node at (68,\yC+2) {12 bits};
  % ram_write
  \draw[thick, opAdc] (0,\yD) -- (93,\yD) -- (93,\yD+4) -- (95,\yD+4) -- (95,\yD) -- (136,\yD);

  % marcas de fase
  \foreach \x/\t in {0/0, 40/96, 44/98, 93/{\(\approx\)146}, 130/{divisor}} {
    \draw[densely dashed, black!35] (\x,\yA+7) -- (\x,\yD-3);
    \node[below] at (\x,\yD-3) {\t};
  }
  \node[below] at (68,\yD-8) {fase (ciclos de 50\,MHz; fora de escala)};

  % anotacoes
  \draw[<-, black!70] (0.5,\yA+2) -- (14,\yA+2) node[right] {instante da amostra};
  \draw[decorate, decoration={brace, amplitude=3pt}] (0,\yA+12) -- (40,\yA+12)
    node[midway, above=2pt] {conversão: 1,92\,$\mu$s};
  \draw[decorate, decoration={brace, amplitude=3pt}] (44,\yB+6) -- (92,\yB+6)
    node[midway, above=2pt] {12 pulsos a 12,5\,MHz};
\end{tikzpicture}
\caption{Um quadro de \texttt{adc\_captura.v} (estado \texttt{S\_CAPT}). Na
fase 0, CONVST sobe e a conversão começa (é o instante da amostra); na fase 96, desce; da fase 98 em
diante, 12 pulsos de SCLK leem o resultado pelo SDO (B11 a B0)
e, nos 6 primeiros, enviam a configuração pelo SDI; no fim da
leitura, a amostra vai para a RAM. Na fase $\text{divisor}-1$ o contador volta
a 0 e o próximo quadro começa. Com divisor 250 (200\,kHz), o quadro de cerca de
146 ciclos cabe com folga.}
\label{fig:quadro-adc}
\end{figure}

\paragraph{O primeiro quadro é descartado.} No LTC2308, a palavra de
configuração enviada durante uma leitura vale para a conversão
\textbf{seguinte}. A conversão do quadro 0 usou a configuração anterior (talvez
de outro canal); por isso são feitas $n+1$ conversões, e a primeira não é
gravada (registrador \texttt{descartar}).

A amostra vai para a RAM crua, nos 12 bits de baixo de uma palavra de 32. O
sinal do modo bipolar é estendido pelo servidor, não aqui.

\paragraph{Modo contínuo.} Com o bit 6 de \texttt{adc\_config} em 1 (usado pela
janela de Aquisição, não pelo analisador), a RAM vira um buffer circular, e o
\texttt{contador} diz quantas amostras já foram gravadas, para o servidor
copiar as novas enquanto a captura segue.

\subsection{O wrapper da FFT: \texttt{fft\_wrapper.v}}
\label{sec:wrapper}

Instanciado em \arq{ghrd_top.v:700}, liga as quatro RAMs (portas B) ao
\textbf{IP FFT da Intel} (\texttt{fft\_core}, gerado no Platform Designer a
partir de \arq{Quartus/fft_core.qsys}). A configuração do IP: 1024 pontos,
arquitetura \emph{Buffered Burst}, entrada de 24 bits, fatores de giro de 18
bits, representação \emph{Block Floating Point}, ordem natural na entrada e na
saída, bidirecional (o mesmo IP faz a IFFT com \texttt{inverse} = 1).

O IP usa a interface \textbf{Avalon-ST}: um fluxo com \texttt{valid} e
\texttt{ready} dos dois lados e marcas de início (\texttt{sop}) e fim
(\texttt{eop}) de pacote. O wrapper converte ``uma RAM endereçável'' em ``um
fluxo de 1024 amostras'' e o contrário, com a máquina de estados da
figura~\ref{fig:fsm-fft}.

\begin{figure}[htbp]
\centering
\begin{tikzpicture}[
    est/.style={caixa, fpga, minimum width=27mm, font=\scriptsize},
    tr/.style={->, thick},
    lab/.style={font=\scriptsize, align=center},
    node distance=9mm and 16mm,
  ]
  \node[est] (idle) {\texttt{S\_IDLE}};
  \node[est, right=of idle] (rs) {\texttt{S\_READ\_START}\\pede \texttt{xn[a]}};
  \node[est, right=of rs] (rw) {\texttt{S\_READ\_WAIT}\\espera a RAM};
  \node[est, below=of rw] (feed) {\texttt{S\_FEED\_FFT}\\\texttt{sink\_valid}, sop/eop};
  \node[est, below=of feed] (recv) {\texttt{S\_RECV\_WAIT}\\\texttt{source\_ready} = 1};
  \node[est, left=of recv] (ws) {\texttt{S\_WRITE\_START}\\escreve \texttt{yn[b]}};
  \node[est, left=of ws] (ww) {\texttt{S\_WRITE\_WAIT}\\espera a RAM};
  \node[est, above=of ww] (comp) {\texttt{S\_COMPLETE}\\\emph{done} = 1};

  \draw[tr] (idle) -- node[above, lab]{\emph{start} $\uparrow$} (rs);
  \draw[tr] (rs) -- (rw);
  \draw[tr] (rw) -- node[right, lab]{lido} (feed);
  \draw[tr] (feed.west) -- ++(-6mm,0) |- node[pos=0.25, left, lab]{\texttt{sink\_ready},\\$a < 1023$:\\$a$++} (rs.south);
  \draw[tr] (feed) -- node[right, lab]{\texttt{sink\_ready}, $a = 1023$} (recv);
  \draw[tr] (recv) -- node[below, lab, yshift=-1mm]{\texttt{source\_valid}} (ws);
  \draw[tr] (ws) -- (ww);
  \draw[tr] (ww.south) -- ++(0,-6mm) -| node[pos=0.25, below, lab]{$b < 1023$: $b$++} (recv.south);
  \draw[tr] (ww) -- node[left, lab]{$b = 1023$} (comp);
  \draw[tr] (comp) -- node[left, lab]{\emph{start} $\downarrow$} (idle);
\end{tikzpicture}
\caption{Máquina de estados de \texttt{fft\_wrapper.v}. A metade de cima
alimenta o IP com as 1024 amostras ($a$ = \texttt{feed\_addr}); a de baixo
recebe os 1024 bins e os grava ($b$ = \texttt{recv\_addr}). No primeiro bin
(\texttt{source\_sop}), o expoente BFP é guardado em \texttt{fft\_exponent}.}
\label{fig:fsm-fft}
\end{figure}

Três detalhes custaram horas de depuração e estão no código por isso:

\begin{itemize}
\item \textbf{\texttt{source\_ready} é um fio combinacional}
  (\texttt{state == S\_RECV\_WAIT}), e não um registrador. Com o \emph{ready}
  oscilando no meio do quadro, a IFFT perdia o fim do pacote, o wrapper ficava
  preso em \texttt{S\_RECV\_WAIT} e toda FFT seguinte dava timeout. A versão
  correta é a do bitstream de 08/09/2026, versionada só em 15/09.
\item \textbf{\texttt{inverse} é travado na borda de subida do \emph{start}.}
  Por isso o servidor escreve esse PIO antes de subir o \emph{start}.
\item \textbf{O IP recebe só os 24 bits de baixo} de cada palavra de entrada, e
  a saída de 24 bits é estendida para 32 antes de ir à RAM. Por isso o cliente
  codifica em Q15.8 com o sinal estendido.
\end{itemize}

% =============================================================================
\section{O caminho numérico completo}
\label{sec:escalas}

Juntando as camadas, um bin $X[k]$ passa pelas escalas da
figura~\ref{fig:escalas}.

\begin{figure}[htbp]
\centering
\begin{tikzpicture}[node distance=3.5mm,
    passo/.style={caixa, minimum width=78mm, font=\small},
    onde/.style={font=\scriptsize, anchor=west}]
  \node[passo, cliente] (p1) {$x_w[n]$: volts, \texttt{float64}};
  \node[passo, cliente, below=of p1] (p2) {$\times\, s \approx 32\,000/\max|x_w|$; $\times 256$; arredonda};
  \node[passo, rede, below=of p2] (p3) {inteiro Q15.8 de 24 bits, pela rede, até a RAM};
  \node[passo, fpga, below=of p3] (p4) {IP FFT: DFT com divisões por 2 contadas em \texttt{exp}};
  \node[passo, hps, below=of p4] (p5) {$\times\, 2^{-\texttt{exp}} / 256$ $\to$ \texttt{float32}};
  \node[passo, cliente, below=of p5] (p6) {$/\, s$ $\to$ $X[k]$ na escala de $x_w$};
  \node[passo, cliente, below=of p6] (p7) {$/\, (\tfrac12 \sum w)$ $\to$ volts de pico na tela};
  \foreach \a/\b in {p1/p2,p2/p3,p3/p4,p4/p5,p5/p6,p6/p7} \draw[seta] (\a) -- (\b);
  \foreach \p/\t in {p1/PC, p2/PC, p3/rede, p4/FPGA, p5/servidor C, p6/PC, p7/PC}
    \node[onde] at ($(\p.east)+(3mm,0)$) {\t};
\end{tikzpicture}
\caption{As escalas de um bin, da janela até a tela. Cada multiplicação feita
na ida é desfeita na volta; o único efeito líquido é o erro de quantização.}
\label{fig:escalas}
\end{figure}

O ADC dá cerca de 9 bits efetivos (ENOB 9,03 e SINAD 56,1\,dB, medidos em
07/10/2026 com senoide de gerador), e a FFT da placa concorda com o NumPy a
cerca de 97\,dB. Portanto \textbf{o limite do analisador é o conversor, não a
FFT}: o erro da FPGA fica uns 40\,dB abaixo do ruído do próprio sinal.

% =============================================================================
\section{Tempo de um quadro}

Medido na placa~1, a partir do PC Windows, CH0, 10\,kHz, 1024 pontos
(30/09/2026):

\begin{table}[H]
\centering
\small
\begin{tabular}{@{}l l@{}}
\toprule
etapa & tempo \\
\midrule
captura no ADC & $1025 / 10\,000$\,Hz $\approx$ 103\,ms, mais a rede \\
FFT (codificação, rede, cálculo, resposta) & 30 a 40\,ms \\
resultado & cerca de 5,6 quadros por segundo \\
\bottomrule
\end{tabular}
\caption{Tempo de um quadro com 1024 pontos a 10\,kHz.}
\end{table}

O pedido de FFT sozinho, com a placa livre, mede de 21 a 25\,ms (testes do
\texttt{morphe-up.sh} no diário). O log do servidor separa as partes de cada
pedido (\texttt{log\_tempos}: \texttt{entrada}, \texttt{fpga},
\texttt{saida}), mas o tempo só do hardware ainda não foi registrado no
diário. A captura, ao contrário, é limitada pela física: $N$ amostras a $f_s$
levam $N/f_s$ segundos. Com 16\,384 pontos a 1\,kHz, um quadro leva 16\,s só
para ser capturado.

% =============================================================================
\section{Limites e cuidados}

\begin{itemize}
\item \textbf{Não há filtro antialiasing} entre o conector J15 e o LTC2308.
  Qualquer componente acima de $f_s/2$ volta como alias, e a tela não tem como
  distingui-lo. A fs deve ficar bem acima do dobro da maior frequência
  presente.
\item \textbf{Faixa de entrada}: 0 a 4,096\,V por pino, em qualquer modo.
  Tensão negativa pode danificar o conversor. Com gerador de funções: offset de
  cerca de 2\,V e no máximo 4\,Vpp, conferidos no osciloscópio antes de ligar;
  terra no pino 10 do J15, CH0 no pino 2.
\item \textbf{Tamanho}: até 32\,768 amostras por captura única (a RAM
  \texttt{adc\_buf}); o analisador usa até 16\,384.
\item \textbf{fs}: de 1\,kHz (divisor 50\,000) a 200\,kHz (divisor 250), sempre
  50\,MHz dividido por um inteiro.
\item \textbf{Entrada sem nada ligado} dá um sinal que começa em cerca de
  2,4\,V e decai nas primeiras amostras (entrada de alta impedância; diário de
  30/09). Não é defeito do analisador.
\item \textbf{Placa compartilhada}: com várias pessoas na mesma placa, os
  pedidos se intercalam, e a taxa de quadros de cada uma cai.
\end{itemize}

% =============================================================================
\section{Para seguir no código}

\begin{table}[H]
\centering
\small
\begin{tabularx}{\linewidth}{@{}X l@{}}
\toprule
quero entender\dots & começar por \\
\midrule
o laço da janela & \arq{Python/analisador_window.py:395} \\
as contas do espectro & \arq{Python/analisador_window.py:120} \\
a fs real e a palavra do ADC & \arq{Python/aquisicao.py:63} e \arq{:86} \\
a normalização Q15.8 & \arq{Python/morphe_protocol.py:198} e \arq{:213} \\
a FFT longa & \arq{Python/blocos.py:241} \\
o despacho no servidor & \arq{C/morphe_server.c:1524} \\
a captura no servidor & \arq{C/morphe_server.c:1138} \\
a FFT no servidor e o BFP & \arq{C/morphe_server.c:1334} \\
o instante de amostragem & \arq{Quartus/adc_captura.v}, estado \texttt{S\_CAPT} \\
a ponte RAM $\leftrightarrow$ IP FFT & \arq{Quartus/fft_wrapper.v} \\
os endereços & \arq{C/hps_0.h} e \arq{C/address_map_arm.h} \\
precisão numérica & \arq{docs/PRECISAO-NUMERICA.md} \\
o ADC em detalhe & \arq{docs/ADC.md} \\
\bottomrule
\end{tabularx}
\caption{Pontos de entrada no código.}
\end{table}

\end{document}
